% GENERATED FILE. Edit canonical lessons, then run npm run generate:books.
% canonical-source-hash: fnv1a64:a2b01334d4ddc8d2
% canonical-lessons: JA-C428-wakuwaku-suru, JA-C428-koufun-suru, JA-C428-kandou-suru, JA-C428-hashagu, JA-C428-harikiru

\chapter{Excited and Moved}
\label{ch:ja-excited-and-moved}
\hlchaptermodality{428}

\begin{chapteropening}
\textbf{By the end of this chapter:} \emph{I can say what excited or moved me.}

Complete the last lesson of chapter 428: I can say what excited or moved me.
\end{chapteropening}

\section[\texorpdfstring{\ja{わくわくする} (wakuwaku suru)}{わくわくする (wakuwaku suru)}]{\ja{わくわくする} (wakuwaku suru) — to be thrilled; to be excited}

\label{lesson:JA-C428-wakuwaku-suru}



\begin{quote}
Before the new one: say the Japanese for to get irritated, then the Japanese for to hesitate.
\end{quote}

\subsection*{You'll want to know: \ja{わくわくする}}
\textbf{\ja{わくわくする}} — \emph{wakuwaku suru} — \textquotedblleft{}to be thrilled; to be excited\textquotedblright{}.

\begin{quote}
\textbf{\ja{あしたの} \ja{りょこうに} \ja{わくわくしています。}} — \emph{ashita no ryokō ni wakuwaku shite imasu} — \textquotedblleft{}I am thrilled about tomorrow's trip.\textquotedblright{}
\end{quote}

Polite \textbf{\ja{わくわくします}}, \emph{wakuwakushimasu}; a \ja{する} verb. Past (polite): \textbf{\ja{こどもが} \ja{わくわくしました。}}, \emph{kodomo ga wakuwaku shimashita}.

\begin{cousinweb}
\textbf{\ja{わくわくする}} (\emph{wakuwaku suru}) is a native Japanese mimetic word, \emph{wakuwaku}, inherited rather than borrowed from Chinese, made a verb with \emph{suru}, \textquotedblleft{}to do\textquotedblright{}; its older origin is not settled (JMdict).
\end{cousinweb}

\begin{grammarlens}[title={What you've built}]
One more word for describing what you lived and felt.
\end{grammarlens}

\subsection*{Your turn}
\begin{itemize}
\raggedright
  \item \emph{Say it:} \emph{wakuwaku suru}
  \item \emph{Say it:} \emph{wakuwaku suru}, once more
  \item \emph{Say it:} \emph{ashita no ryokō ni wakuwaku shite imasu} — I am thrilled about tomorrow's trip.
  \item \emph{Recall:} say the Japanese for to get irritated, then the Japanese for to hesitate, then say \emph{wakuwaku suru} again
\end{itemize}

\begin{tcolorbox}[breakable,colback=teal!4,colframe=teal!35!black,title={Before you move on}]
What does \ja{わくわくする} mean? (\textquotedblleft{}To be thrilled; to be excited\textquotedblright{}.) Say it once more.
\end{tcolorbox}
\section[\texorpdfstring{\ja{こうふんする} (kōfun suru)}{こうふんする (kōfun suru)}]{\ja{こうふんする} (kōfun suru) — to get worked up; to get excited}

\label{lesson:JA-C428-koufun-suru}



\begin{quote}
Before the new one: say the Japanese for to hesitate, then the Japanese for to be thrilled.
\end{quote}

\subsection*{You'll want to know: \ja{こうふんする}}
\textbf{\ja{こうふんする}} — \emph{kōfun suru} — \textquotedblleft{}to get worked up; to get excited\textquotedblright{}.

\begin{quote}
\textbf{\ja{しあいを} \ja{みて、みんな} \ja{こうふんしました。}} — \emph{shiai o mite, minna kōfun shimashita} — \textquotedblleft{}Everyone got worked up at the match.\textquotedblright{}
\end{quote}

Polite \textbf{\ja{こうふんします}}, \emph{kōfunshimasu}; a \ja{する} verb. Past (polite): \textbf{\ja{おとうとが} \ja{こうふんしました。}}, \emph{otōto ga kōfun shimashita}.

\begin{cousinweb}
\textbf{\ja{こうふんする}} (\emph{kōfun suru}) is Sino-Japanese (kango) \emph{kōfun} with the native verb \emph{suru}, \textquotedblleft{}to do\textquotedblright{}: two Chinese characters read with their on'yomi, readings borrowed from Chinese. They are \emph{kō}, \textquotedblleft{}revive\textquotedblright{} (Middle Chinese \emph{xing}) and \emph{fun}, \textquotedblleft{}stirred up\textquotedblright{} (\emph{pjunH}) (KANJIDIC2; Middle Chinese in Baxter's transcription).
\end{cousinweb}

\begin{grammarlens}[title={What you've built}]
One more word for describing what you lived and felt.
\end{grammarlens}

\subsection*{Your turn}
\begin{itemize}
\raggedright
  \item \emph{Say it:} \emph{kōfun suru}
  \item \emph{Say it:} \emph{kōfun suru}, once more
  \item \emph{Say it:} \emph{shiai o mite, minna kōfun shimashita} — Everyone got worked up at the match.
  \item \emph{Recall:} say the Japanese for to hesitate, then the Japanese for to be thrilled, then say \emph{kōfun suru} again
\end{itemize}

\begin{tcolorbox}[breakable,colback=teal!4,colframe=teal!35!black,title={Before you move on}]
What does \ja{こうふんする} mean? (\textquotedblleft{}To get worked up; to get excited\textquotedblright{}.) Say it once more.
\end{tcolorbox}
\section[\texorpdfstring{\ja{かんどうする} (kandō suru)}{かんどうする (kandō suru)}]{\ja{かんどうする} (kandō suru) — to be deeply moved}

\label{lesson:JA-C428-kandou-suru}



\begin{quote}
Before the new one: say the Japanese for to be thrilled, then the Japanese for to get worked up.
\end{quote}

\subsection*{You'll want to know: \ja{かんどうする}}
\textbf{\ja{かんどうする}} — \emph{kandō suru} — \textquotedblleft{}to be deeply moved\textquotedblright{}.

\begin{quote}
\textbf{\ja{あの} \ja{えいがに} \ja{かんどうしました。}} — \emph{ano eiga ni kandō shimashita} — \textquotedblleft{}I was deeply moved by that film.\textquotedblright{}
\end{quote}

Polite \textbf{\ja{かんどうします}}, \emph{kandōshimasu}; a \ja{する} verb. Past (polite): \textbf{\ja{ははが} \ja{かんどうしました。}}, \emph{haha ga kandō shimashita}.

\begin{cousinweb}
\textbf{\ja{かんどうする}} (\emph{kandō suru}) is Sino-Japanese (kango) \emph{kandō} with the native verb \emph{suru}, \textquotedblleft{}to do\textquotedblright{}: two Chinese characters read with their on'yomi, readings borrowed from Chinese. They are \emph{kan}, \textquotedblleft{}emotion\textquotedblright{} (Middle Chinese \emph{komX}) and \emph{dō}, \textquotedblleft{}move\textquotedblright{} (\emph{duwngX}) (KANJIDIC2; Middle Chinese in Baxter's transcription).
\end{cousinweb}

\begin{grammarlens}[title={What you've built}]
One more word for describing what you lived and felt.
\end{grammarlens}

\subsection*{Your turn}
\begin{itemize}
\raggedright
  \item \emph{Say it:} \emph{kandō suru}
  \item \emph{Say it:} \emph{kandō suru}, once more
  \item \emph{Say it:} \emph{ano eiga ni kandō shimashita} — I was deeply moved by that film.
  \item \emph{Recall:} say the Japanese for to be thrilled, then the Japanese for to get worked up, then say \emph{kandō suru} again
\end{itemize}

\begin{tcolorbox}[breakable,colback=teal!4,colframe=teal!35!black,title={Before you move on}]
What does \ja{かんどうする} mean? (\textquotedblleft{}To be deeply moved\textquotedblright{}.) Say it once more.
\end{tcolorbox}
\section[\texorpdfstring{\ja{はしゃぐ} (hashagu)}{はしゃぐ (hashagu)}]{\ja{はしゃぐ} (hashagu) — to frolic; to make merry}

\label{lesson:JA-C428-hashagu}



\begin{quote}
Before the new one: say the Japanese for to get worked up, then the Japanese for to be deeply moved.
\end{quote}

\subsection*{You'll want to know: \ja{はしゃぐ}}
\textbf{\ja{はしゃぐ}} — \emph{hashagu} — \textquotedblleft{}to frolic; to make merry\textquotedblright{}.

\begin{quote}
\textbf{\ja{こどもが} \ja{こうえんで} \ja{はしゃいでいます。}} — \emph{kodomo ga kōen de hashaide imasu} — \textquotedblleft{}The children are frolicking in the park.\textquotedblright{}
\end{quote}

Polite \textbf{\ja{はしゃぎます}}, \emph{hashagimasu}; a godan verb. Past (polite): \textbf{\ja{こどもが} \ja{はしゃぎました。}}, \emph{kodomo ga hashagimashita}.

\begin{cousinweb}
\textbf{\ja{はしゃぐ}} (\emph{hashagu}) is a native Japanese word (wago), inherited rather than borrowed from Chinese; JMdict writes it in kana (JMdict).
\end{cousinweb}

\begin{grammarlens}[title={What you've built}]
One more word for describing what you lived and felt.
\end{grammarlens}

\subsection*{Your turn}
\begin{itemize}
\raggedright
  \item \emph{Say it:} \emph{hashagu}
  \item \emph{Say it:} \emph{hashagu}, once more
  \item \emph{Say it:} \emph{kodomo ga kōen de hashaide imasu} — The children are frolicking in the park.
  \item \emph{Recall:} say the Japanese for to get worked up, then the Japanese for to be deeply moved, then say \emph{hashagu} again
\end{itemize}

\begin{tcolorbox}[breakable,colback=teal!4,colframe=teal!35!black,title={Before you move on}]
What does \ja{はしゃぐ} mean? (\textquotedblleft{}To frolic; to make merry\textquotedblright{}.) Say it once more.
\end{tcolorbox}
\section[\texorpdfstring{\ja{はりきる} (harikiru)}{はりきる (harikiru)}]{\ja{はりきる} (harikiru) — to be full of energy; to be eager}

\label{lesson:JA-C428-harikiru}



\begin{quote}
Before the new one: say the Japanese for to be deeply moved, then the Japanese for to frolic.
\end{quote}

\subsection*{You'll want to know: \ja{はりきる}}
\textbf{\ja{はりきる}} — \emph{harikiru} — \textquotedblleft{}to be full of energy; to be eager\textquotedblright{}.

\begin{quote}
\textbf{\ja{おとうとは} \ja{はりきって} \ja{れんしゅうしています。}} — \emph{otōto wa harikitte renshū shite imasu} — \textquotedblleft{}My younger brother is practising with great energy.\textquotedblright{}
\end{quote}

Polite \textbf{\ja{はりきります}}, \emph{harikirimasu}; a godan verb. Past (polite): \textbf{\ja{せんしゅが} \ja{はりきりました。}}, \emph{senshu ga harikirimashita}.

\begin{cousinweb}
\textbf{\ja{はりきる}} (\emph{harikiru}) is a native Japanese compound (wago), inherited rather than borrowed from Chinese: \emph{hari}, \textquotedblleft{}stretching (haru)\textquotedblright{} and \emph{kiru}, \textquotedblleft{}to cut\textquotedblright{} (JMdict, KANJIDIC2).
\end{cousinweb}

\begin{grammarlens}[title={What you've built}]
One more word for describing what you lived and felt.
\end{grammarlens}

\subsection*{Your turn}
\begin{itemize}
\raggedright
  \item \emph{Say it:} \emph{harikiru}
  \item \emph{Say it:} \emph{harikiru}, once more
  \item \emph{Say it:} \emph{otōto wa harikitte renshū shite imasu} — My younger brother is practising with great energy.
  \item \emph{Recall:} say the Japanese for to be deeply moved, then the Japanese for to frolic, then say \emph{harikiru} again
\end{itemize}

\begin{tcolorbox}[breakable,colback=teal!4,colframe=teal!35!black,title={Before you move on}]
What does \ja{はりきる} mean? (\textquotedblleft{}To be full of energy; to be eager\textquotedblright{}.) Say it once more.
\end{tcolorbox}
